\section{Generalización a otras operaciones binarias}

Ambas estructuras admiten las operaciones propuestas, aunque no siempre
conservan las consultas constantes de RMQ. Esta generalización es teórica:
los experimentos midieron únicamente mínimos.

\paragraph{Propiedades algebraicas.}
Sea \(f:S\times S\to S\) una operación cerrada y
\textbf{asociativa}, \(f(f(a,b),c)=f(a,f(b,c))\). Esta propiedad permite
reagrupar bloques consecutivos: \((S,f)\) es un semigrupo.
Con un \textbf{neutro} \(e\), \(f(e,x)=f(x,e)=x\), se obtiene un
monoide, que facilita inicializar acumuladores y representar rangos vacíos
\cite{atcoder-segtree}. Para rangos no vacíos basta iniciar con el primer bloque.

La \textbf{idempotencia}, \(f(x,x)=x\), permite solapar bloques sin
alterar el resultado. No se necesitan inversos ni conmutatividad si se
preserva el orden. Sin asociatividad, la descomposición habitual puede dar
respuestas incorrectas; reemplazar el mínimo por \(f\) no basta.

\paragraph{Operaciones propuestas.}
MCD y MCM denotan máximo común divisor y mínimo común múltiplo.
Todas las operaciones son asociativas y conmutativas. Sparse Table combina bloques
solapados solo para las idempotentes; en otro caso usa bloques disjuntos
\cite{cp-sparse-table}.

\begin{center}
\begingroup
\small
\setlength{\tabcolsep}{7pt}
\renewcommand{\arraystretch}{1.12}
\begin{tabular}{@{}lccc@{}}
    \toprule
    Operación & Idempotente & Neutro & Consulta Sparse Table \\
    \midrule
    Mínimo & Sí & \(+\infty\) & \(\Theta(1)\) \\
    Máximo & Sí & \(0\) & \(\Theta(1)\) \\
    Suma & No & \(0\) & \(O(\log n)\) \\
    Producto & No & \(1\) & \(O(\log n)\) \\
    XOR & No & \(0\) & \(O(\log n)\) \\
    AND & Sí & \(2^w-1\) & \(\Theta(1)\) \\
    OR & Sí & \(0\) & \(\Theta(1)\) \\
    MCD & Sí & \(0\) & \(\Theta(1)\) \\
    MCM & Sí & \(1\) & \(\Theta(1)\) \\
    \bottomrule
\end{tabular}
\endgroup
\end{center}

Se suponen enteros no negativos de \(w\) bits para AND y OR,
\(\gcd(0,0)=0\) y \(\operatorname{mcm}(0,x)=0\). MCM es idempotente:
\(\operatorname{mcm}(x,x)=x\). Para mínimo, el neutro puede ser el máximo
representable en un dominio acotado, o añadirse \(+\infty\).

\paragraph{Garantías de Sparse Table.}
Construcción, espacio y reconstrucción por cambio conservan
\(\Theta(n\log n)\). Las consultas idempotentes siguen siendo
\(\Theta(1)\). Para suma, producto y XOR, \texttt{query} combina bloques
disjuntos en orden: procesa \(\operatorname{popcount}(k)\) bloques,
con cota \(O(\log(k+1))\subseteq O(\log n)\).
No debe usarse \texttt{query\_idempotent}: sumar los bloques
\([2,3]\) y \([3,5]\) devuelve 13 en vez de 10, por repetir el 3.

\paragraph{Garantías y adaptación de Segment Tree.}
La asociatividad permite conservar construcción y espacio \(\Theta(n)\),
consultas \(O(\log n)\) y actualizaciones \(\Theta(\log n)\)
\cite{cp-segment-tree,atcoder-segtree}. En el código actual, sustituir
\texttt{\_min} requiere corregir \texttt{query}: inicializa con una hoja
que vuelve a incluir, repetición inocua para mínimo pero incorrecta para
suma, producto y XOR. Se debe iniciar con el neutro o el primer bloque
procesado. Para operaciones no conmutativas hacen falta acumuladores
izquierdo y derecho ordenados; las actualizaciones deben combinar hijo
izquierdo y luego derecho, no según cuál fue modificado. Son adaptaciones
propuestas, no cambios a los algoritmos medidos.

\paragraph{Límites del modelo.}
Las cotas suponen agregados y combinaciones de costo constante. Suma,
producto y MCM pueden desbordarse; la aritmética modular no da el resultado
exacto. Dividir antes de multiplicar en MCM no garantiza representabilidad
\cite{cpp-lcm}. Con precisión arbitraria deben contabilizarse el tamaño,
la copia y la combinación de resultados. Al analizar bits, MCD y MCM
también incorporan el costo aritmético \cite{cp-euclid}; las garantías
de RMQ no se trasladan incondicionalmente a enteros ilimitados.
